% !TeX root = ../../main.tex
\subsection{Động lực, mục tiêu và phạm vi của đề tài}

\subsubsection{Động lực của đề tài}

Trên một sàn thương mại điện tử, một đơn hàng không dừng lại ở thao tác mua và trả tiền. Sau khi khách đặt hàng, sàn phải duyệt thanh toán, người bán phải chuẩn bị hàng và giao cho đơn vị vận chuyển, đơn vị vận chuyển đưa hàng tới khách, và cuối cùng khách để lại đánh giá. Mỗi bước được ghi lại ở một nơi khác nhau, do một bên khác nhau phụ trách.

Quy mô của các sàn thương mại điện tử như Shopee, Amazon, Olist\dots đủ lớn để việc theo dõi hoạt động bằng các báo cáo rời rạc không còn phù hợp. Tuy nhiên trong thực tế, việc khai thác dữ liệu để phục vụ công tác quản lý sàn vẫn còn tồn tại nhiều hạn chế. Chẳng hạn, để biết ngành hàng nào vừa hay giao trễ vừa bị đánh giá thấp, người phân tích phải tự nối ít nhất năm bảng: đơn hàng, dòng sản phẩm, sản phẩm, bảng dịch ngành hàng và đánh giá, đồng thời phải tự xử lý trường hợp một đơn có nhiều sản phẩm hoặc nhiều đánh giá. Nếu mỗi người tự viết truy vấn theo cách của mình, cùng một chỉ số có thể ra hai con số khác nhau.

Một khó khăn khác là các chỉ số thường được xem tách rời. Doanh thu nằm ở báo cáo kinh doanh, thời gian giao hàng nằm ở báo cáo vận hành, còn điểm đánh giá nằm ở bộ phận chăm sóc khách hàng. Một vấn đề vận hành ở một tuyến giao hàng có thể kéo theo đánh giá xấu và mất khách ở cả một khu vực, nhưng nếu không đặt các chỉ số cạnh nhau thì rất khó nhìn ra.

Hệ thống giao dịch của sàn lại được thiết kế để ghi nhận trạng thái hiện tại của đơn hàng, không phải để so sánh xu hướng theo tháng, quý hay năm. Dữ liệu nguồn cũng chưa sạch: có giá trị thiếu, dòng trùng, tên thành phố viết nhiều kiểu và mốc thời gian sai thứ tự (Bảng~\ref{tab:chatluong}). Vì vậy cần một kho dữ liệu để gom dữ liệu về một nơi, làm sạch một lần và dùng chung một cách tính chỉ số cho mọi bộ phận.

\subsubsection{Mục tiêu của đề tài}

Mục tiêu chung của đề tài là giúp nhà quản lý sàn theo dõi doanh thu, thời gian giao hàng, mức độ hài lòng của khách hàng và hiệu quả của người bán theo vùng và theo ngành hàng. Các mục tiêu cụ thể gồm:

\begin{enumerate}
\item Tích hợp dữ liệu đơn hàng, sản phẩm, khách hàng, người bán, đánh giá và địa lý vào một kho dữ liệu duy nhất, với cách tính chỉ số thống nhất.
\item Cho phép phân tích doanh thu theo thời gian, khu vực và ngành hàng để thấy được xu hướng tăng trưởng, mùa cao điểm và các ngành hàng chủ lực.
\item Đo được thời gian của từng công đoạn giao hàng và tỷ lệ giao trễ, chi tiết đến bang, tuyến giao hàng và người bán.
\item Theo dõi điểm đánh giá và làm rõ mức độ ảnh hưởng của chất lượng giao hàng lên điểm đánh giá.
\item Xếp hạng người bán bằng một bộ chỉ số chung về doanh thu, giao hàng và đánh giá.
\item Tạo nền tảng dữ liệu cho dashboard và cho thành phần hệ hỗ trợ quyết định (DSS).
\end{enumerate}

\subsubsection{Phạm vi của đề tài}\label{sec:phamvi}
Đề tài tập trung vào hoạt động phân tích sau bán hàng trên một sàn thương mại điện tử, sử dụng bộ dữ liệu có sẵn làm nguồn phân tích. Trong phạm vi này, kho dữ liệu phục vụ bốn mảng phân tích: doanh thu (doanh thu, phí vận chuyển, số đơn và giá trị đơn theo thời gian, khu vực, ngành hàng), giao hàng (thời gian từng công đoạn và tình trạng giao đúng hạn), mức độ hài lòng (điểm đánh giá theo ngành hàng, người bán, khu vực và tình trạng giao hàng) và hiệu quả người bán.

\subsubsection{Các bên liên quan}
Kho dữ liệu không chỉ phục vụ nhu cầu báo cáo của một bộ phận mà còn mang lại giá trị cho nhiều bên liên quan (stakeholder). Mỗi nhóm stakeholder đều nhận được những lợi ích cụ thể từ việc triển khai và sử dụng hệ thống.

Ban điều hành cần một nguồn số liệu thống nhất để theo dõi tình hình kinh doanh và chọn thị trường, ngành hàng trọng điểm. Bộ phận vận hành cần biết tuyến nào, khu vực nào và người bán nào hay giao trễ để điều chỉnh cam kết giao hàng. Bộ phận quản lý ngành hàng và người bán cần căn cứ để hỗ trợ, khuyến khích hoặc cảnh báo người bán. Bộ phận chăm sóc khách hàng cần biết đánh giá thấp tập trung ở đâu để ưu tiên xử lý.

Người bán và khách hàng không trực tiếp dùng hệ thống nhưng vẫn hưởng lợi. Người bán được đánh giá theo tiêu chí rõ ràng và được hỗ trợ khi chất lượng giảm, còn khách hàng nhận được dịch vụ giao hàng tốt hơn khi sàn xử lý được các điểm nghẽn.
